\documentclass[10pt,a4paper]{amsart}
\usepackage[margin=19mm]{geometry}
\usepackage{amsmath,amssymb,mathtools}
\usepackage[scr=rsfs,cal=euler]{mathalfa}
\usepackage{xfrac}
\usepackage[unicode,hidelinks,bookmarks=false]{hyperref}
\newcommand{\ie}{i.e.\ }
\newcommand{\cf}{cf.\ }
\newcommand{\resp}{respectively\ }
\newcommand{\Subsection}{\subsection}
\providecommand{\nobreakdash}{\nobreak\hbox{-}}
\setlength{\emergencystretch}{2em}
\multlinegap0pt
\usepackage{bm}
\usepackage{bbm}
\usepackage{dsfont}
\usepackage{xspace}
\usepackage{cancel}
\usepackage{mathtools}
\usepackage{amsthm}
\makeatletter
\@ifundefined{theorem}{\newtheorem{theorem}{Theorem}}{}
\makeatother
\newcommand{\graph}{\mathcal{G}}
\newcommand{\tti}{\rightarrow\infty}
\title[Entropy of Random Geometric Graphs in High and Low Dimension]{Entropy of Random Geometric Graphs in High and Low Dimensions}
\author{Oliver Baker, Carl P. Dettmann}
\date{Source: arXiv:2503.11418v3 (CC BY 4.0)}
\begin{document}
\maketitle

\noindent\emph{Source and licence.} Entropy of Random Geometric Graphs in High and Low Dimensions. Version arXiv:2503.11418v3 (CC BY 4.0), distributed under the Creative Commons Attribution 4.0 International licence, as recorded in the arXiv licence metadata for this exact version and shown on the abstract page https://arxiv.org/abs/2503.11418v3. Full rights notice: licenses/arxiv-2503.11418-CC-BY-4.0.txt.\par\smallskip

\noindent\textit{Editorial scope.} Counted unit: the Theorem whose source label is \texttt{no convergence for gaussian in torus thm} in the article (source0.tex lines 531--534, source label \texttt{no convergence for gaussian in torus thm}), delivered together with the article's own complete original proof of it (source0.tex lines 535--608, one \texttt{proof} environment of 74 lines). No mathematics is written, repaired or re-derived: statement and proof are the article's own text line for line. Required context retained uncounted: eq:general\_distribution (lines 323--326). Every definition, display and label of those lines is retained unchanged. Omitted from this package and not redistributed: 1--322, 327--530, 609--1547. Nothing inside a retained excerpt is omitted or altered: every retained body is delivered exactly as the article's lines are written, so no omission receipt of any kind accompanies this package. Citations: the retained excerpts carry no citation command at all, so no bibliography entry of the article is transcribed and no reference is redistributed. Reference adaptation: none was needed and none was made; every reference inside the delivered unit resolves inside it, so no unresolved reference or citation is produced. Printed slips kept literal: no printed word, symbol, hypothesis or number of the counted statement or of its proof was altered. Layout and class: the article's own class is \texttt{article} with packages including graphicx, authblk, array, xcolor, geometry, hyperref, tocbibind, amsmath, amsthm, amssymb, amsfonts, breqn, pdflscape, algorithm; the wrapper is the lane's pinned \texttt{amsart} editorial wrapper at 10pt on A4, which restates the theorem-like environments the retained text needs and adds the article's own macro definitions that the retained text needs (\texttt{\textbackslash graph}, \texttt{\textbackslash tti}), with extra wrapper packages (bm, bbm, dsfont, xspace, cancel, mathtools). The delivered numbering is the unit's own. Assets: no retained span contains a figure environment, a graphics-inclusion command, a PDF-page inclusion, a table or a TikZ picture; no image file is copied into this package, the package declares no asset dependency and no image file is redistributed with it. Licence: version arXiv:2503.11418v3 of this article is distributed under the Creative Commons Attribution 4.0 International licence, as recorded in the arXiv licence metadata for this exact version and shown on the abstract page https://arxiv.org/abs/2503.11418v3. A full rights notice is delivered at licenses/arxiv-2503.11418-CC-BY-4.0.txt. Characters: every retained excerpt is pure ASCII, so the delivered engine, XeLaTeX with its default Latin Modern text font, reports no missing character.
\begin{equation}
    \label{eq:general_distribution}
    \nu(\vec{x}) = \prod_{i=1}^d \pi(x_i)
\end{equation}

\begin{theorem}
\label{no convergence for gaussian in torus thm}
    Suppose the ensemble $\graph$ of $n$ node hard RGGs on the $\mathbb{T}^d$ has nodes distributed according to a distribution of the form (\ref{eq:general_distribution}). Then $\graph$ converges in distribution to $G(n,p)$ as $d\tti$ if and only if the distribution of nodes in each coordinate is uniform.
\end{theorem}

\begin{proof}
    We have convergence to the ER ensemble if and only if the diagonal entries of the covariance matrix $\beta = 0$. Recall that, if $X,Y,Z$ are i.i.d. according to $\pi$,
    \begin{equation}
        \beta = \mathbb{E}[(\rho_T(X,Y)^2-\mu_T)(\rho_T(X,Y)^2-\mu_T)]
    \end{equation}
    So, by writing the expectation as an integral, and expanding the brackets, we have convergence to ER if and only if
    \begin{equation}
            \int_0^1\int_0^1\int_0^1\rho_T(x,y)^2\rho_T(x,z)^2 d\pi(x)\pi(y)\pi(z)dxdydz = \mu_T^2
    \end{equation}
    \begin{equation}
        \iff \int_0^1\left(\int_0^1 \rho_T(x,y)^2\pi(y)dy\right)^2\pi(x)dx = \mu_T^2
    \end{equation}
    \begin{equation}
        \iff \mathbb{E}_X[\mathbb{E}_Y[\rho_T(X,Y)^2]^2] = \mathbb{E}_X[\mathbb{E}_Y[\rho_T(X,Y)^2]]^2
    \end{equation}
    which is equivalent to
    \begin{equation}
        Var_X(\mathbb{E}_Y[\rho_T(X,Y)^2]) = 0
    \end{equation}
    and so $\mathbb{E}_Y[\rho_T(X,Y)^2]$ is almost surely constant. Rewriting this as an integral,
    \begin{equation}
        \int_0^1 \rho_T(x,y)^2\pi(y)dy = C
    \end{equation}
    $\pi$ almost-everywhere. In fact, this constant must be $\mu_T$, since integrating this function (which is constant) over the distribution of $x$ must give $\mu_T$. To show that the unique distribution satisfying this property is uniform, we will calculate the Fourier series of both sides. For any constant $C$, we have the $\xi^{\text{th}}$ Fourier mode of $C$, $\hat{C}(\xi)$ is:
    \begin{equation}
        \hat{C}(\xi) = \int_0^1 Ce^{-2\pi i \xi x}dx = \frac{C}{2\pi i \xi} \left(1-e^{-2\pi i\xi}\right) = 0
    \end{equation}
    whenever $\xi$ is a non-zero integer, and $\hat{C}_0 = C$. For the left hand side, write
    \begin{equation}
        F(x) = \int_0^1\rho_T(x,y)^2\pi(y)dy = \int_0^1 \rho_T(x-y)^2\pi(y)dy = (\rho_T^2 \ast \pi)(x) 
    \end{equation}
    where in the second inequality we have used a slight abuse of notation to show that $\rho_T(x,y)$ is dependent only on the difference between $x$ and $y$, and $\ast$ denotes the convolution of two functions. Here, we treat the functions $\pi$ and $\rho_T$ as periodic on $\mathbb{R}$ so we can apply the convolution theorem, but restrict them to $[0,1]$ so that we are working on $\mathbb{T}$. We have:
    \begin{equation}
        \label{eq:fourier_product}
        \hat{F}(\xi) = \widehat{\rho_T^2}(\xi)\hat{\pi}(\xi)
    \end{equation}
    where for a function $f$ we denote by $\hat{f}(\xi)$ its $\xi^{\text{th}}$ Fourier mode. We can write the Fourier modes of $\rho_T^2$ by splitting it over the cases when the argument is greater than, or less than $\frac{1}{2}$:
    \begin{equation}
        \widehat{\rho_T^2}(\xi) = \int_0^1 \rho_T(x)^2e^{-2\pi i \xi x}dx = \int_0^{\frac{1}{2}} |x|^2e^{-2\pi i \xi x}dx + \int_{\frac{1}{2}}^1(1-|x|)^2e^{-2\pi i \xi  x}dx
    \end{equation}
    \begin{equation}
        = \int_{-\frac{1}{2}}^{\frac{1}{2}} x^2e^{-2 \pi i \xi x}dz
    \end{equation}
    \begin{equation}
        = \begin{cases}
            \frac{1}{12} & \xi = 0 \\
            \frac{(-1)^{\xi}}{2\pi^2\xi} & \xi \in \mathbb{Z}\setminus\{0\}
        \end{cases}
    \end{equation}
    % which is $\frac{1}{12}$ when $\xi = 0$, and on computing the integral using standard techniques,
    % \begin{equation}
    %     \widehat{\rho_T^2}(\xi ) = -2\frac{e^{-\pi i \xi}}{(-2\pi i \xi)^2}
    % \end{equation}
    % if $\xi \neq0$. Since $\xi $ is an integer, we can reduce this to
    % \begin{equation}
    %     \widehat{\rho_T^2}(\xi ) = \frac{(-1)^{\xi }}{2\pi^2\xi ^2} 
    % \end{equation}
    which is non-zero for every $\xi  \in \mathbb{Z}\setminus{0}$. Therefore (\ref{eq:fourier_product}) implies
    \begin{equation}
        \frac{1}{12}\hat{\pi}(0) = \mu_T
    \end{equation}
    and
    \begin{equation}
        \frac{(-1)^{\xi }}{2\pi^2\xi ^2}  \hat{\pi}(\xi ) = 0
    \end{equation}
    for $\xi \neq 0$. Therefore, we must have
    \begin{equation}
        \hat{\pi}(\xi)  = \begin{cases}
            12\mu_T & \text{if }\xi =0, \\
            0 & \text{if }\xi \neq0
        \end{cases}
    \end{equation}
    and therefore $\pi$ must be constant on $[0,1]$. Calculating $\mu_T$ for the uniform distribution, we find $\mu_T = \frac{1}{12}$, and therefore $\hat{\pi}(\xi) = 1$ if $\xi = 0$, and is $0$ elsewhere, so $\pi(x) = 1$ on $[0,1]$.
\end{proof}

\end{document}
